The most important file in this folder, because the thing that stops people is not lack of
knowledge. It is making forty loops and finishing nothing.

Assumes three to five hours a week, in sessions of an hour or more. If you get less than
that in a given week, cut the reading, not the making.

\begin{reviewbox}
\textbf{In short.} You finish and export four tracks in eight weeks. Finished means arranged,
mixed badly, exported to MP3, and filed. Not good. Finished.
\end{reviewbox}

\subsection{Week 1: get all the way through the pipeline once}\label{pr01_eight_weeks-week-1-get-all-the-way-through-the-pipeline-once}

Read \texttt{FL\_STUDIO/00\_setup.md} and do the setup, including the template.

Do \texttt{FL\_STUDIO/02\_first\_beat.md} end to end.

Do \texttt{FL\_STUDIO/03\_first\_track.md} end to end, including the export. Give yourself one
session for this and do not stop halfway.

Read \texttt{THEORY/01\_what\_edm\_is.md}.

\textbf{Deliverable: track 1, exported.} It will be bad. File it in a folder called \texttt{01\_finished}
and do not touch it again.

The point of week 1 is to prove the pipeline works. Every subsequent week improves one part
of a process you have already completed once.

\subsection{Week 2: drums}\label{pr01_eight_weeks-week-2-drums}

Read \texttt{THEORY/03\_rhythm.md} properly.

Make \textbf{five different eight-bar drum loops}. Not tracks, just loops. Different tempos:
118, 122, 124, 126, 128. Save them all.

Do the ear-training exercise from \texttt{exercises.md}: clap the offbeat through three full tracks.

Then pick one Chicago house record and rebuild its drum pattern by ear, using only the
bundled 909 samples. Spend at least an hour on this. It is the single highest-value exercise
in the whole plan.

\textbf{Deliverable: five loops and one transcription.}

\subsection{Week 3: harmony}\label{pr01_eight_weeks-week-3-harmony}

Read \texttt{THEORY/04\_harmony.md}.

In FL, make a two-bar chord loop in A minor using seventh chords, and voice-lead between
them properly. Then transpose the same loop to F minor and to C minor without changing the
shape.

Write five different basslines over the same chord loop: root notes on offbeats, octave
bounce, a walking line, one that syncopates against the kick, and one that plays only on
beat 1.

Then combine week 2\textquotesingle s best drum loop with this week\textquotesingle s best chords and bass.

\textbf{Deliverable: track 2, arranged and exported.} Use the arrangement table from
\texttt{FL\_STUDIO/03\_first\_track.md} again. Reusing a structure is not cheating.

\subsection{Week 4: sound design}\label{pr01_eight_weeks-week-4-sound-design}

Read \texttt{THEORY/05\_sound\_design.md}.

Do the one-hour 3x Osc exercise at the end of that file. All six steps, slowly, listening.

Then spend a session with \textbf{Transistor Bass} making an acid line. Sixteenth notes, high
resonance, and automate the cutoff across sixteen bars. This is the Chicago instrument and
it deserves its own session.

Make three sounds from scratch with no presets: a bass, a pad, a pluck.

\textbf{Deliverable: three patches and one acid loop.} No finished track this week.

\subsection{Week 5: arrangement}\label{pr01_eight_weeks-week-5-arrangement}

Read \texttt{THEORY/06\_arrangement.md}.

Take your best loop so far and arrange it using the \textbf{top down} method: copy the full loop
across 128 bars, then delete elements to carve out the intro, breakdown and build.

Focus this week entirely on the build and the drop. Spend an hour just on bars 81 to 97.
Riser, filter automation, hat roll, silence before the drop, crash on the downbeat.

\textbf{Deliverable: track 3, exported.} This one should be noticeably better than track 1.

\subsection{Week 6: mixing}\label{pr01_eight_weeks-week-6-mixing}

Read \texttt{THEORY/07\_mixing.md}.

Take track 3 and remix it from scratch levels. Highpass everything except kick and bass. Set
up two reverb sends. Get the sidechain working properly with Fruity Limiter rather than
volume automation.

Do the referencing exercise: load a commercial Chicago house track into your project,
level-match, and A/B for a full hour.

\textbf{Deliverable: track 3, remixed.} Compare the two exports back to back.

\subsection{Week 7: speed}\label{pr01_eight_weeks-week-7-speed}

Read \texttt{THEORY/08\_genre\_map.md} and pick a subgenre you have not tried.

Then do the constraint exercise: \textbf{make a complete track in one three-hour sitting.} One
synth, four drum sounds, no presets browsing, no going back. Export whatever you have at the
three hour mark.

Speed exercises work because they remove the option of perfectionism. You will be surprised
how much of your normal session is spent deciding rather than doing.

\textbf{Deliverable: track 4, exported.}

\subsection{Week 8: consolidate}\label{pr01_eight_weeks-week-8-consolidate}

Listen to tracks 1 through 4 in order. Write down, for each, one thing that is better than
the previous one and one thing that is still wrong.

That list is your curriculum for the next eight weeks.

Then pick whichever of the four has the most in it, and spend the whole week making it
actually good. This is the first time in the plan you are allowed to polish.

Read \texttt{LISTENING/chicago\_scene\_today.md} and go to one thing. A club night, a record shop, a
festival.

\textbf{Deliverable: one track you are not embarrassed by, and a plan for what to learn next.}

\subsection{Rules that make this work}\label{pr01_eight_weeks-rules-that-make-this-work}

\textbf{Finish badly rather than not at all.} A finished bad track teaches you more than an
unfinished good loop, because the last 20\% is where all the skills you are missing live.

\textbf{Time-box.} Set a timer. When a session ends, save and stop, even mid-idea. Especially
mid-idea, because you will want to come back.

\textbf{Do not restart projects.} The urge to start fresh because "this one is not working" is
the enemy. Finish it, then start fresh.

\textbf{Keep a \texttt{01\_finished} folder} and put everything in it. Being able to hear month one
against month six is the main reward available in the early period, and it is a real one.

\textbf{Do not buy anything.} Not plugins, not sample packs, not a controller. You have more than
you can use. Buying things feels like progress and is not.

\textbf{One reference track per session.} Pick a Chicago record, have it loaded, A/B against it.

\subsection{What eight weeks will and will not get you}\label{pr01_eight_weeks-what-eight-weeks-will-and-will-not-get-you}

It will get you: fluency in FL Studio, a working mental model of how house is built, four
finished tracks, and the ability to hear what is happening in a record.

It will not get you: tracks that sound like the records you like. That takes years, and
everyone underestimates it. The gap between your taste and your ability is at its widest
right now and it closes slowly.

That gap is the normal condition of learning anything. You have a PhD; you already know what
the first year of not being good at something feels like.
